%=============================================================================

配分函数（partition function）是统计物理中最核心的概念之一。它是一个数学对象，包含了系统的所有热力学信息。一旦我们知道了配分函数，就可以通过微分运算得到所有热力学量，如内能、熵、自由能等。在本节中，我们将从经典统计力学出发，逐步介绍配分函数的概念，然后推广到量子统计力学，最后介绍路径积分（path integral）表述，这是连接统计物理与量子场论（quantum field theory）的桥梁。

\section{统计力学的基本概念}

统计力学是连接微观物理学和宏观热力学的桥梁。它的基本思想是：宏观系统的热力学性质是大量微观粒子集体行为的统计平均结果。

\begin{definition}[微观状态与宏观状态]
\textbf{微观状态}（microstate）描述系统中每个粒子的详细状态。对于经典系统，微观状态由所有粒子的位置坐标 $q_i$ 和动量 $p_i$ 给出；对于量子系统，微观状态由系统的量子态 $|\psi\rangle$ 描述。

\textbf{宏观状态}（macrostate）描述系统的宏观可观测量，如温度 $T$、压强 $P$、体积 $V$、内能 $E$ 等。一个宏观状态通常对应大量微观状态。
\end{definition}

\begin{definition}[热力学极限]
当粒子数 $N \to \infty$ 和体积 $V \to \infty$，但粒子数密度 $N/V$ 保持有限时的极限称为\textbf{热力学极限}（thermodynamic limit）。在热力学极限下，涨落可以忽略不计，热力学描述变得精确。
\end{definition}

\begin{note}
统计力学的基本假设是\textbf{等概率原理}（equal a priori probability principle）：对于一个孤立系统，所有满足宏观约束的微观状态出现的概率相等。这个假设是统计力学的基石，虽然它不能从更基本的原理推导出来，但它的推论与实验高度一致。
\end{note}

\section{统计系综}

统计系综（statistical ensemble）是描述系统概率分布的数学框架。它是大量相同系统的集合，每个系统都处于相同的宏观条件下，但可能处于不同的微观状态。图 \ref{fig:ensembles} 展示了三种主要统计系综的示意图。

\begin{figure}[htbp]
\centering
\begin{subfigure}[b]{0.3\textwidth}
\centering
\begin{tikzpicture}[scale=0.8]
    \draw[thick] (0,0) rectangle (2,2);
    \node at (1,1) {系统};
    \draw[thick] (0,0.5) -- (-0.5,0.5);
    \draw[thick] (0,1.5) -- (-0.5,1.5);
    \node at (-0.8,1) {绝热壁};
    \node at (1,-0.3) {微正则系综};
    \node at (1,-0.6) {$E, V, N$ 固定};
\end{tikzpicture}
\end{subfigure}
\hfill
\begin{subfigure}[b]{0.3\textwidth}
\centering
\begin{tikzpicture}[scale=0.8]
    \draw[thick] (0,0) rectangle (2,2);
    \node at (1,1) {系统};
    \draw[thick, red] (0,0.5) -- (-0.8,0.5);
    \draw[thick, red] (0,1.5) -- (-0.8,1.5);
    \node at (-1.1,1) {热库 $T$};
    \node at (1,-0.3) {正则系综};
    \node at (1,-0.6) {$T, V, N$ 固定};
\end{tikzpicture}
\end{subfigure}
\hfill
\begin{subfigure}[b]{0.3\textwidth}
\centering
\begin{tikzpicture}[scale=0.8]
    \draw[thick] (0,0) rectangle (2,2);
    \node at (1,1) {系统};
    \draw[thick, red] (0,0.5) -- (-0.8,0.5);
    \draw[thick, red] (0,1.5) -- (-0.8,1.5);
    \node at (-1.1,1) {热库 $T$};
    \draw[thick, blue] (2,0.5) -- (2.8,0.5);
    \draw[thick, blue] (2,1.5) -- (2.8,1.5);
    \node at (3.1,1) {粒子库 $\mu$};
    \node at (1,-0.3) {巨正则系综};
    \node at (1,-0.6) {$T, V, \mu$ 固定};
\end{tikzpicture}
\end{subfigure}
\caption{三种主要统计系综的示意图。从左到右：微正则系综（孤立系统）、正则系综（与热库接触）、巨正则系综（与热库和粒子库接触）。}
\label{fig:ensembles}
\end{figure}

\begin{definition}[微正则系综]
\textbf{微正则系综}（microcanonical ensemble）描述孤立系统，能量 $E$ 固定。系统处于每个微观状态的概率相等（等概率原理）：
\begin{equation}
    \rho = \frac{1}{\Omega(E)} \delta(H - E),
\end{equation}
其中 $\Omega(E)$ 是能量为 $E$ 的微观状态数，$\rho$ 是概率密度函数。熵定义为：
\begin{equation}
    S = k_B \ln \Omega(E),
\end{equation}
这就是著名的玻尔兹曼熵公式。$k_B$ 是玻尔兹曼常数（Boltzmann constant），$k_B \approx 1.381 \times 10^{-23}$ J/K。
\end{definition}

\begin{definition}[正则系综]
\textbf{正则系综}（canonical ensemble）描述与热库接触的系统，温度 $T$ 固定。系统处于能量为 $E_i$ 的微观状态的概率为：
\begin{equation}
    p_i = \frac{e^{-\beta E_i}}{Z},
\end{equation}
其中 $\beta = 1/(k_B T)$ 是逆温度（inverse temperature），$Z$ 是配分函数（partition function）：
\begin{equation}
    Z = \sum_i e^{-\beta E_i}.
\end{equation}
这个概率分布称为玻尔兹曼分布（Boltzmann distribution）。
\end{definition}

\begin{definition}[巨正则系综]
\textbf{巨正则系综}（grand canonical ensemble）描述与粒子库和热库接触的系统，温度 $T$ 和化学势 $\mu$ 固定。系统处于粒子数为 $N$、能量为 $E_i$ 的微观状态的概率为：
\begin{equation}
    p_i = \frac{e^{-\beta (E_i - \mu N)}}{\mathcal{Z}},
\end{equation}
其中 $\mathcal{Z}$ 是巨配分函数（grand partition function）：
\begin{equation}
    \mathcal{Z} = \sum_{N=0}^{\infty} \sum_i e^{-\beta (E_i - \mu N)} = \sum_{N=0}^{\infty} z^N Z_N,
\end{equation}
其中 $z = e^{\beta \mu}$ 是逸度（fugacity），$Z_N$ 是粒子数为 $N$ 时的正则配分函数。
\end{definition}

\begin{property}[系综等价性]
\textbf{系综等价性}（ensemble equivalence）：在热力学极限下，不同系综给出相同的热力学性质。这意味着我们可以根据方便选择任何一个系综来计算。
\end{property}

\begin{remark}
系综等价性的物理原因在于：在热力学极限下，相对涨落按 $1/\sqrt{N}$ 趋于零。因此，固定粒子数 $N$（正则系综）与允许粒子数涨落（巨正则系综）的差别可以忽略。但在有限系统或临界点附近，系综等价性可能失效，此时必须谨慎选择系综。
\end{remark}

\section{经典统计力学}

考虑一个由 $N$ 个粒子组成的经典系统。每个粒子的位置坐标为 $q_i$，动量为 $p_i$（其中 $i = 1, 2, \ldots, N$）。系统的总能量由哈密顿量（Hamiltonian）$H(q, p)$ 给出，它是所有粒子的坐标和动量的函数。

\begin{definition}[相空间]
\textbf{相空间}（phase space）：由所有粒子的位置和动量张成的 $6N$ 维空间。相空间中的一个点代表系统的一个微观状态。
\end{definition}

\begin{definition}[哈密顿量]
\textbf{哈密顿量}（Hamiltonian）：系统的总能量，通常是动能和势能之和：
\begin{equation}
    H(q, p) = \sum_{i=1}^{N} \frac{p_i^2}{2m} + V(q_1, \ldots, q_N),
\end{equation}
其中第一项是动能，$V$ 是势能。
\end{definition}

在正则系综中，系统与一个温度为 $T$ 的热库（heat bath）接触。根据玻尔兹曼分布（Boltzmann distribution），系统处于某个微观状态 $(q, p)$ 的概率与 $e^{-\beta H(q, p)}$ 成正比。

\begin{theorem}[经典正则配分函数]
\textbf{配分函数}定义为所有可能微观状态的玻尔兹曼因子之和。对于经典系统，这个"求和"实际上是相空间（phase space）上的积分：
\begin{equation}
    Z = \int \frac{d^N q \, d^N p}{h^{3N} N!} e^{-\beta H(q, p)}.
    \label{eq:classical_Z}
\end{equation}
其中：
\begin{itemize}
    \item $d^N q \, d^N p$ 表示对整个相空间积分，$d^N q = dq_1 \cdots dq_N$，$d^N p = dp_1 \cdots dp_N$
    \item $h$ 是普朗克常数（Planck constant），$h \approx 6.626 \times 10^{-34}$ J$\cdot$s。除以 $h^{3N}$ 是为了使配分函数无量纲化，这来自于量子力学中相空间体积元的最小单位
    \item $N!$ 是全同粒子的修正，这是因为交换两个全同粒子不产生新的微观状态。如果不除以 $N!$，会导致熵不是广延量，这就是著名的吉布斯佯谬（Gibbs paradox）
\end{itemize}
\end{theorem}

\begin{note}
配分函数 $Z$ 是所有可能微观状态的"统计权重"的总和。能量越低的状态，其玻尔兹曼因子 $e^{-\beta E}$ 越大，因此对配分函数的贡献越大。温度越低（$\beta$ 越大），低能状态的相对权重越大。
\end{note}

\begin{property}[非相互作用系统的配分函数分解]
如果系统的哈密顿量可以分解为独立部分之和 $H = H_1 + H_2 + \cdots + H_N$（即各部分之间无相互作用），则配分函数分解为各部分配分函数的乘积：
\begin{equation}
    Z = Z_1 \cdot Z_2 \cdots Z_N = \prod_{i=1}^N Z_i,
\end{equation}
其中 $Z_i = \int \frac{dq_i \, dp_i}{h} e^{-\beta H_i}$。相应地，自由能具有可加性：$F = F_1 + F_2 + \cdots + F_N$。这一性质大大简化了无相互作用系统（如理想气体）的计算。
\end{property}

\begin{example}[经典理想气体的配分函数]
对于 $N$ 个无相互作用的经典粒子，哈密顿量仅为动能之和。由于动能与位置无关，配分函数可以分解为动量部分和位置部分：
\begin{equation}
    Z = \frac{1}{N! h^{3N}} \left(\int d^3q \right)^N \left(\int d^3p \, e^{-\beta p^2/(2m)}\right)^N = \frac{V^N}{N!} \left(\frac{2\pi m}{\beta h^2}\right)^{3N/2}.
\end{equation}
这展示了配分函数如何将微观力学（哈密顿量）与宏观热力学（通过 $\beta$）联系起来。
\end{example}

\section{热力学量的推导}

\begin{theorem}[热力学量与配分函数]
一旦知道了配分函数，所有热力学量都可以从它推导出来。配分函数之所以如此重要，是因为它是一个"生成函数"：通过对它进行适当的微分运算，我们可以得到所有感兴趣的物理量。

\textbf{亥姆霍兹自由能}（Helmholtz free energy）$F$：
\begin{equation}
    F = -k_B T \ln Z = -\frac{1}{\beta} \ln Z.
\end{equation}
自由能的物理意义是：在等温等容过程中，系统对外做的最大功等于自由能的减少。自由能是热力学势（thermodynamic potential），它以温度 $T$ 和体积 $V$ 为自然变量。

从自由能出发，我们可以得到其他热力学量：

\textbf{熵}（entropy）$S$：
\begin{equation}
    S = -\left(\frac{\partial F}{\partial T}\right)_V = k_B (\ln Z + \beta E).
\end{equation}
熵是系统无序程度的度量。根据热力学第二定律，孤立系统的熵永不减少。

\textbf{内能}（internal energy）$E$：
\begin{equation}
    E = -\frac{\partial}{\partial \beta} \ln Z = \langle H \rangle = \sum_i p_i E_i.
\end{equation}
内能是系统能量的统计平均值。

\textbf{压强}（pressure）$P$：
\begin{equation}
    P = -\left(\frac{\partial F}{\partial V}\right)_T = \frac{1}{\beta} \frac{\partial \ln Z}{\partial V}.
\end{equation}
压强是系统对体积变化的响应。

\textbf{比热容}（heat capacity）$C_V$：
\begin{equation}
    C_V = \left(\frac{\partial E}{\partial T}\right)_V = k_B \beta^2 \left( \langle H^2 \rangle - \langle H \rangle^2 \right) = k_B \beta^2 \langle (\Delta H)^2 \rangle.
\end{equation}
比热容度量系统吸收热量的能力，它与能量涨落成正比。

\textbf{化学势}（chemical potential）$\mu$：
\begin{equation}
    \mu = -T \left(\frac{\partial S}{\partial N}\right)_{E,V} = \left(\frac{\partial F}{\partial N}\right)_{T,V}.
\end{equation}
化学势度量增加一个粒子时系统自由能的变化。
\end{theorem}

\begin{proof}[热力学量从配分函数推导的证明]
我们证明内能 $E = -\partial \ln Z / \partial \beta$。由配分函数的定义 $Z = \sum_i e^{-\beta E_i}$，对 $\beta$ 求导：
\begin{equation}
    \frac{\partial Z}{\partial \beta} = \sum_i (-E_i) e^{-\beta E_i},
\end{equation}
因此
\begin{equation}
    \frac{\partial \ln Z}{\partial \beta} = \frac{1}{Z} \frac{\partial Z}{\partial \beta} = -\frac{1}{Z} \sum_i E_i e^{-\beta E_i} = -\sum_i E_i p_i = -\langle H \rangle = -E.
\end{equation}
这就证明了 $E = -\partial \ln Z / \partial \beta$。

接下来证明自由能 $F = -k_B T \ln Z$ 与热力学的关系。在正则系综中，概率分布 $p_i = e^{-\beta E_i}/Z$ 最大化了信息熵 $S = -k_B \sum_i p_i \ln p_i$，约束条件为 $\sum_i p_i = 1$ 和 $\sum_i p_i E_i = E$。利用拉格朗日乘子法，可以证明自由能 $F = E - TS$ 满足 $F = -k_B T \ln Z$。具体地，将 $p_i = e^{-\beta E_i}/Z$ 代入熵的表达式：
\begin{equation}
    S = -k_B \sum_i p_i (-\beta E_i - \ln Z) = k_B \beta E + k_B \ln Z = \frac{E - F}{T},
\end{equation}
其中最后一步使用了 $F = -k_B T \ln Z$，从而 $F = E - TS$，这正是亥姆霍兹自由能的标准定义。

比热容的公式可以通过对内能求导得到：
\begin{equation}
    C_V = \frac{\partial E}{\partial T} = \frac{\partial}{\partial T}\left(-\frac{\partial \ln Z}{\partial \beta}\right) = k_B \beta^2 \frac{\partial^2 \ln Z}{\partial \beta^2}.
\end{equation}
利用 $\partial^2 \ln Z / \partial \beta^2 = \langle H^2 \rangle - \langle H \rangle^2$，即得 $C_V = k_B \beta^2 \langle (\Delta H)^2 \rangle$。
\end{proof}

\section{经典例子}

\begin{example}[理想气体]
\textbf{例子 1：理想气体}（ideal gas）。对于 $N$ 个无相互作用的经典粒子，哈密顿量为
\begin{equation}
    H(q, p) = \sum_{i=1}^{N} \frac{p_i^2}{2m}.
\end{equation}
配分函数为：
\begin{equation}
    Z = \frac{1}{N!} \left(\frac{V}{h^3}\right)^N \left(\int d^3 p \, e^{-\beta p^2/(2m)}\right)^N = \frac{1}{N!} \left(\frac{V}{\lambda^3}\right)^N,
\end{equation}
其中 $\lambda = h/\sqrt{2\pi m k_B T}$ 是热德布罗意波长（thermal de Broglie wavelength），它度量量子效应的重要性。

从配分函数可以推导出理想气体的状态方程：
\begin{equation}
    PV = Nk_B T,
\end{equation}
以及内能 $E = \frac{3}{2}Nk_B T$（能量均分定理的结果）和熵（Sackur-Tetrode 公式）：
\begin{equation}
    S = Nk_B \left[ \ln\left(\frac{V}{N\lambda^3}\right) + \frac{5}{2} \right].
\end{equation}
Sackur-Tetrode 公式给出了理想气体的绝对熵，这是经典统计力学的重要成就。
\end{example}

\begin{example}[经典谐振子]
\textbf{例子 2：经典谐振子}（classical harmonic oscillator）。一维经典谐振子的哈密顿量为：
\begin{equation}
    H(q, p) = \frac{p^2}{2m} + \frac{1}{2}m\omega^2 q^2,
\end{equation}
其中 $\omega$ 是角频率。配分函数为：
\begin{equation}
    Z = \frac{1}{h} \int dq \, dp \, e^{-\beta H(q, p)} = \frac{1}{\beta \hbar \omega},
\end{equation}
其中 $\hbar = h/(2\pi)$ 是约化普朗克常数。平均能量为 $E = k_B T$，这就是能量均分定理（equipartition theorem）：每个二次自由度贡献 $k_B T/2$ 的能量。
\end{example}

\begin{example}[二能级系统]
\textbf{例子 3：二能级系统}（two-level system）。考虑一个具有两个能级 $E_0 = 0$ 和 $E_1 = \epsilon$ 的系统。配分函数为：
\begin{equation}
    Z = 1 + e^{-\beta \epsilon}.
\end{equation}
平均能量为：
\begin{equation}
    E = \frac{\epsilon e^{-\beta \epsilon}}{1 + e^{-\beta \epsilon}} = \frac{\epsilon}{e^{\beta \epsilon} + 1}.
\end{equation}
在高温极限（$\beta \epsilon \ll 1$）下，$E \approx \epsilon/2$，两个能级等概率占据；在低温极限（$\beta \epsilon \gg 1$）下，$E \approx \epsilon e^{-\beta \epsilon} \to 0$，系统几乎完全占据基态。
\end{example}

\begin{example}[顺磁性物质]
\textbf{例子 4：顺磁性物质}（paramagnetic substance）。考虑 $N$ 个自旋为 $1/2$ 的粒子置于外磁场 $B$ 中，每个粒子有两个能级 $E_\pm = \mp \mu B$。系统的配分函数为：
\begin{equation}
    Z = \left(2\cosh(\beta\mu B)\right)^N.
\end{equation}
由此可得磁化强度 $M = N\mu\tanh(\beta\mu B)$，在高温下得到居里定律 $M \approx N\mu^2 B/(k_B T)$，即磁化率 $\chi \propto 1/T$。
\end{example}

\section{涨落理论}

在统计力学中，物理量围绕其平均值涨落。涨落的大小可以通过配分函数来计算。

\begin{theorem}[能量涨落与比热容]
\textbf{能量涨落}（energy fluctuation）：
\begin{equation}
    \langle (\Delta E)^2 \rangle = \langle H^2 \rangle - \langle H \rangle^2 = \frac{\partial^2 \ln Z}{\partial \beta^2} = k_B T^2 C_V.
\end{equation}
这个公式表明，能量涨落与比热容成正比。
\end{theorem}

\textbf{相对涨落}（relative fluctuation）：
\begin{equation}
    \frac{\sqrt{\langle (\Delta E)^2 \rangle}}{E} \sim \frac{1}{\sqrt{N}},
\end{equation}
对于宏观系统（$N \sim 10^{23}$），相对涨落可以忽略不计（约为 $10^{-12}$），这就是为什么热力学可以描述宏观系统。

\begin{theorem}[涨落-耗散定理]
\textbf{涨落-耗散定理}（fluctuation-dissipation theorem）：系统对外界扰动的响应（如比热容、磁化率等）与系统内部的涨落成正比。这是统计力学的一个基本定理，它将宏观响应函数与微观涨落联系起来。
\end{theorem}

\textbf{磁化率}（magnetic susceptibility）：对于磁性系统，磁化率 $\chi$ 与磁矩涨落成正比：
\begin{equation}
    \chi = \beta \langle (\Delta M)^2 \rangle = \beta (\langle M^2 \rangle - \langle M \rangle^2).
\end{equation}

\begin{remark}[涨落与相变]
涨落在相变（phase transition）附近起着关键作用\cite{Landau1937}。在临界点附近，关联长度趋于无穷，涨落变得宏观可见，相对涨落不再按 $1/\sqrt{N}$ 趋于零。此时热力学描述失效，必须使用重正化群（renormalization group）方法\cite{Wilson1974}来处理临界现象。
\end{remark}

\section{量子统计力学}

在量子力学中，系统的状态由希尔伯特空间（Hilbert space）中的态矢量 $|\psi\rangle$ 描述，可观测量由厄米算符（Hermitian operator）描述。系统的能量由哈密顿算符（Hamiltonian operator）$\hat{H}$ 的本征值给出：
\begin{equation}
    \hat{H} |n\rangle = E_n |n\rangle,
\end{equation}
其中 $|n\rangle$ 是能量本征态（energy eigenstate），$E_n$ 是对应的能量本征值（energy eigenvalue）。

\begin{definition}[密度矩阵]
\textbf{密度矩阵}（density matrix）：在量子统计力学中，系统的状态由密度矩阵 $\hat{\rho}$ 描述。对于正则系综：
\begin{equation}
    \hat{\rho} = \frac{e^{-\beta \hat{H}}}{Z},
\end{equation}
其中 $Z = \Tr e^{-\beta \hat{H}}$ 是配分函数。
\end{definition}

\begin{theorem}[量子正则配分函数]
\textbf{量子配分函数}定义为玻尔兹曼算符（Boltzmann operator）$e^{-\beta \hat{H}}$ 的迹（trace）：
\begin{equation}
    Z = \Tr e^{-\beta \hat{H}} = \sum_n \langle n | e^{-\beta \hat{H}} | n \rangle = \sum_n e^{-\beta E_n}.
\end{equation}
这里我们使用了能量本征态的完备性关系（completeness relation）$\sum_n |n\rangle \langle n| = \hat{1}$。

\textbf{迹}（trace）的定义：对于任意算符 $\hat{A}$，其迹定义为 $\Tr \hat{A} = \sum_n \langle n | \hat{A} | n \rangle$，其中求和遍及完备基。迹的值不依赖于基的选择。
\end{theorem}

量子配分函数是所有能量本征态的玻尔兹曼因子之和。与经典情况类似，低能态对配分函数的贡献更大。

量子配分函数同样包含了系统的所有热力学信息。自由能、熵、内能等热力学量的表达式与经典情况完全相同：
\begin{align}
    F &= -k_B T \ln Z, \\
    S &= -\left(\frac{\partial F}{\partial T}\right)_V, \\
    E &= -\frac{\partial}{\partial \beta} \ln Z = \frac{\Tr(\hat{H} e^{-\beta \hat{H}})}{\Tr e^{-\beta \hat{H}}} = \sum_n p_n E_n,
\end{align}
其中 $p_n = e^{-\beta E_n}/Z$ 是系统处于能量本征态 $|n\rangle$ 的概率。

\begin{property}[量子配分函数的经典极限]
当温度足够高（$\beta \to 0$）时，量子配分函数趋于经典配分函数。具体地，对于由 $N$ 个粒子组成的系统，量子配分函数在高温极限下满足：
\begin{equation}
    Z_{\text{quantum}} \xrightarrow{\beta \to 0} Z_{\text{classical}} = \frac{1}{N! h^{3N}} \int d^Nq \, d^Np \, e^{-\beta H(q,p)}.
\end{equation}
物理上，这是因为高温下热德布罗意波长 $\lambda = h/\sqrt{2\pi m k_B T}$ 远小于粒子间距，量子统计效应（全同粒子的对称化或反对称化要求）可以忽略，系统行为趋于经典\cite{Landau1937}。
\end{property}

\begin{remark}[量子统计与经典统计的区别]
量子统计力学与经典统计力学的根本区别在于：量子力学中，全同粒子不可区分，且态空间是离散的。对于玻色子，波函数在粒子交换下对称；对于费米子，波函数反对称。这些量子统计效应在低温下变得重要，导致玻色-爱因斯坦凝聚和费米简并等现象。
\end{remark}

\section{量子例子}

\begin{example}[量子谐振子]
\textbf{例子 1：量子谐振子}（quantum harmonic oscillator）。一维量子谐振子的能级为：
\begin{equation}
    E_n = \hbar \omega \left(n + \frac{1}{2}\right), \quad n = 0, 1, 2, \ldots
\end{equation}
其中 $\hbar$ 是约化普朗克常数，$\omega$ 是角频率。注意，最低能级 $E_0 = \hbar\omega/2$ 不为零，这称为零点能（zero-point energy），是量子力学的特征。

配分函数为：
\begin{equation}
    Z = \sum_{n=0}^{\infty} e^{-\beta \hbar \omega (n + 1/2)} = \frac{e^{-\beta \hbar \omega/2}}{1 - e^{-\beta \hbar \omega}} = \frac{1}{2 \sinh(\beta \hbar \omega/2)}.
\end{equation}
平均能量为：
\begin{equation}
    E = \frac{\hbar \omega}{2} + \frac{\hbar \omega}{e^{\beta \hbar \omega} - 1}.
\end{equation}
第一项是零点能，第二项是热激发的贡献。在高温极限（$\beta \hbar \omega \ll 1$）下，$E \approx k_B T$，与经典结果一致；在低温极限（$\beta \hbar \omega \gg 1$）下，$E \approx \hbar\omega/2$，系统冻结在基态。
\end{example}

\begin{example}[量子二能级系统]
\textbf{例子 2：量子二能级系统}。考虑自旋 1/2 粒子在磁场 $B$ 中，能级为 $E_\pm = \pm \mu B$，其中 $\mu$ 是磁矩。配分函数为：
\begin{equation}
    Z = e^{\beta \mu B} + e^{-\beta \mu B} = 2\cosh(\beta \mu B).
\end{equation}
磁化强度为：
\begin{equation}
    M = \frac{1}{\beta} \frac{\partial \ln Z}{\partial B} = \mu \tanh(\beta \mu B).
\end{equation}
在高温极限（$\beta \mu B \ll 1$）下，$M \approx \beta \mu^2 B$，即居里定律（Curie's law）；在低温极限（$\beta \mu B \gg 1$）下，$M \to \mu$，即磁化饱和。
\end{example}

\begin{example}[玻色子和费米子]
\textbf{例子 3：玻色子和费米子}。对于量子理想气体，配分函数取决于粒子的统计性质。

\textbf{玻色子}（bosons）：服从玻色-爱因斯坦统计（Bose-Einstein statistics），自旋为整数。每个单粒子态可以被任意多个玻色子占据。配分函数为：
\begin{equation}
    \ln Z = -\sum_k \ln(1 - e^{-\beta(\epsilon_k - \mu)}),
\end{equation}
其中 $\epsilon_k$ 是单粒子能级，$\mu$ 是化学势。平均占据数为玻色-爱因斯坦分布：
\begin{equation}
    \langle n_k \rangle = \frac{1}{e^{\beta(\epsilon_k - \mu)} - 1}.
\end{equation}

\textbf{费米子}（fermions）：服从费米-狄拉克统计（Fermi-Dirac statistics），自旋为半整数。每个单粒子态最多被一个费米子占据（泡利不相容原理）。配分函数为：
\begin{equation}
    \ln Z = \sum_k \ln(1 + e^{-\beta(\epsilon_k - \mu)}).
\end{equation}
平均占据数为费米-狄拉克分布：
\begin{equation}
    \langle n_k \rangle = \frac{1}{e^{\beta(\epsilon_k - \mu)} + 1}.
\end{equation}
\end{example}

\begin{example}[玻色-爱因斯坦凝聚]
\textbf{例子 4：玻色-爱因斯坦凝聚}（Bose-Einstein condensation，BEC）。考虑一个由 $N$ 个无相互作用玻色子组成的系统，处于体积为 $V$ 的盒子中。当温度降低到临界温度 $T_c$ 以下时，宏观量的粒子会凝聚到能量最低的单粒子态（基态），这就是玻色-爱因斯坦凝聚。图 \ref{fig:bec} 展示了 BEC 的能级占据示意图。

\begin{figure}[htbp]
\centering
\begin{tikzpicture}[scale=0.9]
    % 能级
    \draw[thick] (0,0) -- (3,0);
    \node[left] at (0,0) {$E_0$};
    \draw[thick] (0,1) -- (3,1);
    \node[left] at (0,1) {$E_1$};
    \draw[thick] (0,2) -- (3,2);
    \node[left] at (0,2) {$E_2$};
    \draw[thick, dashed] (0,3) -- (3,3);
    \node[left] at (0,3) {$E_3$};
    \node at (1.5,3.5) {$\vdots$};
    % 高温占据
    \fill[blue!30] (0.5,0) rectangle (1.0,0.05);
    \fill[blue!30] (0.5,1) rectangle (0.8,1.05);
    \fill[blue!30] (0.5,2) rectangle (0.7,2.05);
    \node[below] at (0.75,-0.1) {\small 各能级均匀占据};
    % 低温占据
    \fill[red!50] (2.0,0) rectangle (3.0,0.3);
    \fill[red!30] (2.0,1) rectangle (2.3,1.05);
    \fill[red!20] (2.0,2) rectangle (2.15,2.05);
    \node[below] at (2.5,-0.1) {\small 宏观占据基态};
    % 箭头
    \draw[->, thick] (1.5,0.5) -- (1.7,0.5);
    \node[above] at (1.6,0.5) {$T \downarrow$};
    % BEC 标注
    \node[above, red] at (2.5,0.35) {BEC};
    \node[above] at (0.75,0.1) {$T > T_c$};
    \node[above] at (2.5,0.1) {$T < T_c$};
    % 分隔线
    \draw[thin, gray, dashed] (1.6,-0.3) -- (1.6,3.8);
\end{tikzpicture}
\caption{玻色-爱因斯坦凝聚的能级占据示意图。左：$T > T_c$ 时，粒子按玻色-爱因斯坦分布在各能级上；右：$T < T_c$ 时，宏观量粒子凝聚到基态 $E_0$。}
\label{fig:bec}
\end{figure}

临界温度为：
\begin{equation}
    T_c = \frac{2\pi \hbar^2}{m k_B}\left(\frac{N}{V \zeta(3/2)}\right)^{2/3},
\end{equation}
其中 $\zeta(3/2) \approx 2.612$ 是黎曼 $\zeta$ 函数的值。在 $T < T_c$ 时，基态粒子数为：
\begin{equation}
    N_0 = N\left[1 - \left(\frac{T}{T_c}\right)^{3/2}\right].
\end{equation}
BEC 是量子统计力学最深刻的预言之一，1995 年在超冷原子气体中首次实验观测到。

\begin{note}[BEC 的物理意义]
玻色-爱因斯坦凝聚的发生可以通过热德布罗意波长来理解。当 $\lambda \sim n^{-1/3}$（即粒子间距与波长可比）时，不同粒子的波函数开始重叠，量子统计效应变得重要。此时粒子"失去个体性"，形成宏观量子态。BEC 在超导（库珀对凝聚）和超流现象中起着核心作用。
\end{note}
\end{example}

\section{路径积分表述}

路径积分（path integral）是量子力学的一种等价表述，由费曼（Feynman）于 1948 年提出。它的核心思想是：一个粒子从初态到末态的跃迁振幅（transition amplitude）等于所有可能路径的贡献之和。

\begin{definition}[传播子]
\textbf{传播子}（propagator）：粒子从位置 $q_i$ 在时间 $t_i$ 传播到位置 $q_f$ 在时间 $t_f$ 的振幅为：
\begin{equation}
    K(q_f, t_f; q_i, t_i) = \langle q_f | e^{-i\hat{H}(t_f - t_i)/\hbar} | q_i \rangle = \int \mathcal{D}q(t) \, e^{iS[q]/\hbar},
\end{equation}
其中 $S[q] = \int_{t_i}^{t_f} dt \, L(q, \dot{q})$ 是经典作用量，$L$ 是拉格朗日量（Lagrangian），积分是对所有满足边界条件 $q(t_i) = q_i$、$q(t_f) = q_f$ 的路径进行的。
\end{definition}

\begin{note}
\textbf{路径积分测度}（path integral measure）$\mathcal{D}q(t)$：这是对所有可能路径的"求和"，可以理解为将时间分成 $N$ 个小段，对每个时刻的位置积分，然后取 $N \to \infty$ 的极限。具体地，将时间区间 $[t_i, t_f]$ 等分为 $N$ 段，每段长度 $\epsilon = (t_f - t_i)/N$，则
\begin{equation}
    \int \mathcal{D}q(t) = \lim_{N\to\infty} \left(\frac{m}{2\pi i \hbar \epsilon}\right)^{N/2} \prod_{k=1}^{N-1} \int dq_k.
\end{equation}
\end{note}

在统计物理中，路径积分提供了一个深刻的联系：量子统计力学的配分函数可以写成路径积分的形式。这个联系是通过\textbf{虚时}（imaginary time）$\tau = it$ 实现的，这称为 Wick 转动（Wick rotation）。

\begin{theorem}[Wick 转动]
设系统的哈密顿量为 $\hat{H}$，则配分函数 $Z = \Tr e^{-\beta\hat{H}}$ 可以通过 Wick 转动 $t \to -i\tau$ 从量子力学的时间演化算符得到。
\end{theorem}

\begin{proof}[Wick 转动的证明]
考虑量子力学中的时间演化算符 $U(t) = e^{-i\hat{H}t/\hbar}$。在位置基下，传播子为：
\begin{equation}
    K(q_f, t; q_i, 0) = \langle q_f | e^{-i\hat{H}t/\hbar} | q_i \rangle.
\end{equation}
现在做解析延拓（analytic continuation），令 $t = -i\hbar\beta$，即 $t/\hbar = -i\beta$：
\begin{equation}
    K(q_f, -i\hbar\beta; q_i, 0) = \langle q_f | e^{-i\hat{H}(-i\hbar\beta)/\hbar} | q_i \rangle = \langle q_f | e^{-\beta\hat{H}} | q_i \rangle.
\end{equation}
配分函数是迹，即对 $q_f = q_i$ 的情况积分：
\begin{equation}
    Z = \Tr e^{-\beta\hat{H}} = \int dq \, \langle q | e^{-\beta\hat{H}} | q \rangle.
\end{equation}
在路径积分表示中，$e^{-\beta\hat{H}}$ 对应于虚时演化 $t \to -i\hbar\beta$。此时，闵可夫斯基作用量中的 $e^{iS/\hbar}$ 变为：
\begin{equation}
    e^{iS/\hbar} = \exp\left(\frac{i}{\hbar}\int_0^t dt' \, L\right) \xrightarrow{t = -i\hbar\beta} \exp\left(-\int_0^{\beta} d\tau \, L_E\right) = e^{-S_E},
\end{equation}
其中 $S_E = \int_0^{\beta} d\tau \, L_E$ 是欧几里得作用量，$L_E$ 是通过对 $t \to -i\tau$ 的替换从 $L$ 得到的。对于非相对论粒子 $L = \frac{1}{2}m\dot{q}^2 - V(q)$，Wick 转动将动能项 $\frac{1}{2}m(dq/dt)^2$ 变为 $-\frac{1}{2}m(dq/d\tau)^2$，再乘以 $i$ 后变为正的 $\frac{1}{2}m(dq/d\tau)^2$，势能项 $-V$ 变为 $+V$，因此 $L_E = \frac{1}{2}m\dot{q}^2 + V(q)$。
\end{proof}

\begin{remark}[Wick 转动与有限温度]
Wick 转动 $t \to -i\tau$ 将量子力学的实时传播子与统计力学的配分函数联系起来。这个对应关系有一个深刻的物理含义：有限温度 $T = 1/(k_B\beta)$ 的量子系统等价于一个在虚时方向上具有周期 $\beta$ 的欧几里得场论。温度越高（$\beta$ 越小），虚时圆的周长越小，系统的量子特性越弱；温度越低（$\beta$ 越大），虚时圆越大，量子效应越显著。在 AdS/CFT 对偶中，这个对应关系推广为全角（AdS）中的黑洞视界半径与边界 CFT 的温度之间的关系\cite{Hawking1983}。
\end{remark}

\begin{theorem}[配分函数的路径积分表示]
考虑一个量子系统，其哈密顿量为 $\hat{H}$。配分函数 $Z = \Tr e^{-\beta \hat{H}}$ 可以写成路径积分的形式：
\begin{equation}
    Z = \int \mathcal{D}\phi \, e^{-S_E[\phi]},
    \label{eq:path_integral_Z}
\end{equation}
其中 $S_E[\phi]$ 是欧几里得作用量（Euclidean action），积分是对所有在虚时方向上周期为 $\beta$ 的场构型进行的。
\end{theorem}

\begin{proof}[路径积分表示配分函数的证明]
我们从 $Z = \int dq \, \langle q | e^{-\beta\hat{H}} | q \rangle$ 出发。将虚时区间 $[0, \beta]$ 分成 $N$ 段，每段 $\epsilon = \beta/N$：
\begin{equation}
    e^{-\beta\hat{H}} = \left(e^{-\epsilon\hat{H}}\right)^N.
\end{equation}
在每一段中插入完备性关系 $\int dq_k |q_k\rangle\langle q_k| = \hat{1}$：
\begin{equation}
    Z = \int dq_0 \cdots dq_{N-1} \prod_{k=0}^{N-1} \langle q_{k+1} | e^{-\epsilon\hat{H}} | q_k \rangle,
\end{equation}
其中 $q_N = q_0$（周期性边界条件，来自迹运算）。对于小 $\epsilon$，利用 Trotter 公式：
\begin{equation}
    \langle q_{k+1} | e^{-\epsilon\hat{H}} | q_k \rangle \approx \langle q_{k+1} | e^{-\epsilon\hat{T}} e^{-\epsilon\hat{V}} | q_k \rangle,
\end{equation}
其中 $\hat{T} = \hat{p}^2/(2m)$ 是动能算符，$\hat{V} = V(\hat{q})$ 是势能算符。计算矩阵元：
\begin{equation}
    \langle q_{k+1} | e^{-\epsilon\hat{T}} e^{-\epsilon\hat{V}} | q_k \rangle = \sqrt{\frac{m}{2\pi\hbar^2\epsilon}} \exp\left(-\epsilon\left[\frac{m}{2}\left(\frac{q_{k+1}-q_k}{\epsilon}\right)^2 + V(q_k)\right]\right).
\end{equation}
取 $N \to \infty$ 极限，指数上的求和变为积分：
\begin{equation}
    \epsilon \sum_{k=0}^{N-1} \left[\frac{m}{2}\left(\frac{q_{k+1}-q_k}{\epsilon}\right)^2 + V(q_k)\right] \to \int_0^{\beta} d\tau \left[\frac{1}{2}m\dot{q}^2 + V(q)\right] = S_E[q].
\end{equation}
因此配分函数变为：
\begin{equation}
    Z = \int \mathcal{D}q(\tau) \, e^{-S_E[q]},
\end{equation}
其中路径积分测度 $\mathcal{D}q(\tau)$ 包含了所有归一化因子，路径 $q(\tau)$ 满足周期性边界条件 $q(0) = q(\beta)$。
\end{proof}

\begin{definition}[欧几里得作用量]
\textbf{什么是欧几里得作用量？}在通常的量子力学中，作用量 $S$ 是拉格朗日量（Lagrangian）$L$ 对时间的积分：
\begin{equation}
    S = \int dt \, L(q, \dot{q}).
\end{equation}
当我们做 Wick 转动 $t \to -i\tau$ 时，闵可夫斯基作用量（Minkowski action）$S$ 变成欧几里得作用量 $S_E$：
\begin{equation}
    S_E = \int_0^{\beta} d\tau \, L_E(q, \dot{q}),
\end{equation}
其中 $L_E$ 是欧几里得拉格朗日量。对于非相对论粒子，$L_E = \frac{1}{2}m\dot{q}^2 + V(q)$。注意，欧几里得作用量中动能项的符号是正的，这与闵可夫斯基作用量不同。
\end{definition}

\begin{definition}[周期性边界条件]
\textbf{周期性边界条件}（periodic boundary condition）：在配分函数的路径积分中，场构型必须满足周期性边界条件 $\phi(0) = \phi(\beta)$。这是因为迹运算要求初态和末态相同：
\begin{equation}
    Z = \Tr e^{-\beta \hat{H}} = \sum_n \langle n | e^{-\beta \hat{H}} | n \rangle.
\end{equation}
在路径积分表示中，这对应于在虚时方向上周期为 $\beta$ 的场构型。
\end{definition}

\begin{figure}[htbp]
\centering
\begin{tikzpicture}[scale=1.0]
    % 虚时圆
    \draw[thick, ->] (0,0) circle (1.5);
    \node at (0,-0.3) {$\beta$};
    \draw[thick, blue] (45:1.5) -- (45:1.5 |- 0,0) node[below, black] {$\phi$};
    \node[above right] at (45:1.6) {$\tau$};
    % 周期性路径
    \draw[thick, red, domain=0:360, samples=100] plot ({1.5*cos(\x) + 0.3*sin(3*\x)*cos(\x)}, {1.5*sin(\x) + 0.3*sin(3*\x)*sin(\x)});
    \node[right, red] at (1.8,0.5) {$\phi(0)=\phi(\beta)$};
    % 标注
    \node at (0,-2.3) {虚时方向上的周期性场构型};
\end{tikzpicture}
\caption{配分函数路径积分中的周期性边界条件示意图。虚时方向形成一个周长为 $\beta = 1/(k_BT)$ 的圆，场构型在此方向上必须满足周期性条件 $\phi(0) = \phi(\beta)$。}
\label{fig:periodic}
\end{figure}

路径积分表述告诉我们，量子统计力学等价于一个在虚时方向上具有周期性边界条件的经典场论。虚时方向的长度 $\beta = 1/(k_B T)$ 与温度成反比。温度越高，虚时方向越短；温度越低，虚时方向越长。

\section{相图与临界现象}

统计力学的一个重要应用是描述物质的相变行为\cite{Landau1937}。图 \ref{fig:phase_diagram} 展示了典型的物质相图。

\begin{figure}[htbp]
\centering
\begin{tikzpicture}[scale=1.2]
    % 坐标轴
    \draw[thick, ->] (0,0) -- (4.5,0) node[right] {$T$};
    \draw[thick, ->] (0,0) -- (0,4.5) node[above] {$P$};
    % 相界线
    \draw[thick, red, domain=0.5:3.5, samples=50] plot (\x, {3.5/(\x+0.2)-0.5}) node[right] {气液相界};
    \draw[thick, blue] (1.5,0) -- (1.5,3.5) node[above] {固液相界};
    % 区域标注
    \node at (0.8,1) {固};
    \node at (2.5,0.8) {液};
    \node at (3.5,1.5) {气};
    % 关键点
    \fill[red] (1.5,2.8) circle (2pt) node[above right] {三相点};
    \fill[blue] (3.0,2.2) circle (2pt) node[above right] {临界点};
    % 临界点标注
    \draw[thick, dashed] (3.0,0) -- (3.0,2.2);
    \node[below] at (3.0,0) {$T_c$};
    % 临界指数标注
    \draw[<->, thick] (3.0,2.5) -- (3.8,2.5) node[right] {$|T-T_c|^{-\nu}$};
\end{tikzpicture}
\caption{典型的物质相图（$P$-$T$ 图）。气液相界线在临界点 $(T_c, P_c)$ 终结，此处相变变为连续的。在临界点附近，关联长度趋于无穷，系统表现出标度不变性和普适性\cite{Wilson1974}。}
\label{fig:phase_diagram}
\end{figure}

\begin{remark}[配分函数与相变]
相变对应于配分函数（或其导数）的奇异性。对于有限系统，配分函数是解析函数，不存在真正的相变。只有在热力学极限下，$\ln Z$ 的奇异性才可能出现。Lee-Yang 定理严格证明了这一点：相变来自配分函数零点在复平面上逼近实轴。
\end{remark}

\begin{remark}[临界指数与普适性]
在临界点附近，热力学量表现出幂律行为，例如比热容 $C \sim |T - T_c|^{-\alpha}$、序参量 $\phi \sim |T - T_c|^{\beta}$、关联长度 $\xi \sim |T - T_c|^{-\nu}$ 等。令人惊奇的是，临界指数 $\alpha, \beta, \nu, \ldots$ 仅依赖于系统的对称性和空间维度，而与微观细节无关。这一\textbf{普适性}（universality）由重正化群理论给出了解释\cite{Wilson1974}：不同微观系统在重正化群流下流向相同的不动点，因此表现出相同的临界行为。
\end{remark}

\begin{remark}[全息原理与配分函数]
在 AdS/CFT 对偶\cite{Maldacena1997}中，边界共形场论（CFT）的配分函数等于体引力理论的配分函数：
\begin{equation}
    Z_{\text{CFT}} = Z_{\text{gravity}}.
\end{equation}
这就是 GKPW（Gubser-Klebanov-Polyakov-Witten）关系。具体地，如果 CFT 中某个算符 $\mathcal{O}$ 的源为 $\phi_0$，则：
\begin{equation}
    Z_{\text{CFT}}[\phi_0] = \left. Z_{\text{gravity}}[\phi]
\right|_{\phi \to \phi_0 \text{ at boundary}},
\end{equation}
即引力侧的路径积分在边界上取值为 CFT 的源。

特别地，对于有限温度的 CFT，其配分函数对应于 AdS 空间中的黑洞背景\cite{Hawking1983}。黑洞的霍金温度恰好等于 CFT 的温度，黑洞的贝肯斯坦-霍金熵等于 CFT 的热力学熵。这个对偶是理解量子引力和强耦合场论的最有力工具之一。
\end{remark}

\begin{example}[Schwarzschild-AdS 黑洞与 CFT 热力学]
考虑 AdS$_5$ 中的 Schwarzschild 黑洞，其视界半径 $r_h$ 与边界 CFT 的温度之间的关系为 $T = r_h/(\pi L^2)$，其中 $L$ 是 AdS 半径。黑洞的贝肯斯坦-霍金熵为：
\begin{equation}
    S_{\text{BH}} = \frac{\text{Area}}{4G_5} = \frac{\pi^2}{2G_5} V_3 r_h^3,
\end{equation}
其中 $V_3$ 是边界空间的体积。通过 AdS/CFT 对偶，这等于边界 $\mathcal{N}=4$ 超对称 Yang-Mills 理论在强耦合极限下的热力学熵。这验证了全息原理中配分函数相等的深刻含义。
\end{example}

\begin{note}[路径积分在量子场论中的推广]
在量子场论中，配分函数的路径积分表示推广为对所有场构型的泛函积分。对于标量场 $\phi(x)$：
\begin{equation}
    Z[J] = \int \mathcal{D}\phi \, \exp\left(-S_E[\phi] + \int d^dx \, J(x)\phi(x)\right),
\end{equation}
其中 $J(x)$ 是外源场。$Z[J]$ 是量子场论的生成泛函（generating functional），所有关联函数都可以通过对 $J$ 求泛函导数得到。在 AdS/CFT 中，这个生成泛函的对偶引力描述由 GKPW 规则给出\cite{Maldacena1997}。
\end{note}

\section{小结}

在本节中，我们介绍了配分函数在统计物理中的核心地位：
\begin{itemize}
    \item 统计力学是连接微观物理学和宏观热力学的桥梁
    \item 统计系综是描述系统概率分布的数学框架，包括微正则系综、正则系综和巨正则系综
    \item 经典配分函数是相空间上的积分，包含了所有热力学信息
    \item 量子配分函数是玻尔兹曼算符的迹
    \item 涨落理论将宏观响应函数与微观涨落联系起来，在相变附近涨落变得重要\cite{Landau1937}
    \item 路径积分表述将配分函数写成对所有场构型的积分，通过 Wick 转动连接了量子力学和统计力学
    \item 在 AdS/CFT 对偶中，配分函数是连接边界场论与体引力理论的核心对象\cite{Maldacena1997}
\end{itemize}
配分函数之所以如此重要，是因为它是一个"生成函数"：通过对它进行适当的微分运算，我们可以得到所有感兴趣的物理量。在下一节中，我们将看到配分函数在量子场论中扮演着类似的角色。

\begin{note}[配分函数的统一视角]
从更抽象的角度看，配分函数是一个关于外源 $J$ 的泛函 $Z[J]$：
\begin{itemize}
    \item 在统计力学中，$J = \beta$（逆温度），$Z[\beta] = \Tr e^{-\beta\hat{H}}$
    \item 在量子场论中，$J$ 是外源场，$Z[J] = \int \mathcal{D}\phi \, e^{-S[\phi] + J\cdot\phi}$
    \item 在 AdS/CFT 中，$J$ 是边界上算符的源，$Z_{\text{CFT}}[J] = Z_{\text{gravity}}[\phi_0 = J]$
\end{itemize}
这三种情况的共同结构是：配分函数编码了系统对外部扰动的全部线性和非线性响应信息。所有关联函数、响应函数和热力学量都可以通过对 $J$ 的泛函微分得到。这种统一视角是理解量子场论、统计物理和量子引力之间深层联系的关键。
\end{note}

%=============================================================================
